Human motion data—temporal sequences of three-dimensional body-joint
positions—offers a rich yet compact signal for predictive
analytics~\citep{ionescu2013human3}. Unlike raw video,
it encodes fine-grained spatiotemporal structure while remaining orders of
magnitude sparser, making it a natural fit for inference pipelines 
deployable in embodied and real-time systems. This combination has driven adoption across safety-critical
domains. In surveillance, gait-based recognition systems identify
individuals from body movement alone~\citep{sepas2022deep}.
In clinical practice, deep models trained on gait accurately characterize
the severity of Parkinson's disease, enabling remote and continuous
assessment~\citep{eskofier2016recent,el2020deep}.

\paragraph{The need for explainability.}
In these high-stakes settings, predictive accuracy alone is insufficient.
Clinicians, for example, must understand \emph{which} body parts
move anomalously, \emph{when} during a sequence, and \emph{why} that
drove a model's decision --to build trust and improve model performance. Indeed, deep models
are prone to shortcut learning~\citep{geirhos2020shortcut}, focusing onto
spurious correlations that can amplify systemic biases.
For example, \citet{obermeyer2019dissecting} demonstrated
that a widely deployed health algorithm systematically underestimated
illness in Black patients relative to equally sick white patients, cutting
their access to additional care by more than half. Trust in
safety-critical models learning from human motion data therefore requires reliable explanations.

\paragraph{SHAP and its limitations for human motion.}
Explainable AI (XAI) provides a principled toolkit for auditing model
decisions~\citep{arrieta2020explainable,guidotti2018survey}.  Among
feature-attribution approaches, SHAP~\citep{lundberg2017unified} has
become the dominant framework, grounded in cooperative game theory and
supported by well-studied axiomatic
guarantees. Yet, SHAP is fundamentally
ill-suited to human motion.  Its core mechanism---perturbing individual
input features and measuring the resulting change in model
output---violates two structural properties that are intrinsic to human
movement, described next. \textit{Manifold constraint.}  Human motion occupies a low-dimensional
manifold: not every trajectory of 3-D joint coordinates corresponds to
physically plausible movement.  Standard SHAP perturbations sample freely
from the full feature space, generating out-of-distribution inputs that are
never observed in practice. \textit{Temporal coherence.}  Motion is inherently sequential: the
relevance of a joint's trajectory is inseparable from the phase of the
movement in which it occurs. Identifying, for example, that the hip flexion
during mid-stance is diagnostic is fundamentally different from asking
whether any single joint coordinate at any single frame
matters.

Prior work has addressed each challenge in isolation. Manifold-aware methods confine perturbations to the data manifold~\citep{frye2020shapley,taufiq2023manifold,olsen2022using}, while temporal extensions enforce sequential dependencies across time-series inputs~\citep{bento2021timeshap,tempel2025explaining,NAYEBI2023104438}. However, no existing method satisfies both requirements simultaneously, precluding principled attribution for human motion.

\paragraph{The missing benchmark.}
Compounding this problem, the community currently lacks the tools to
\emph{quantify} how badly existing XAI methods fail on human motion
data.  Measuring explanation quality requires access to ground-truth
attributions---a notoriously hard problem.
The XAI literature has addressed this through carefully constructed
synthetic datasets in which the ground-truth important features are
known by design~\citep{arras2022clevr,hedstrom2023quantus,Yang2019BenchmarkingAM}.
No such benchmark exists for human motion, leaving practitioners without
a reliable way to compare or validate attribution methods in this domain.

\paragraph{Contributions.}
We address both gaps with \textsc{MotionSHAP}, a SHAP framework that jointly enforces manifold plausibility and temporal coherence, and pair it with a suite of purpose-built synthetic benchmarks and evaluations on established real-world motion datasets. Our contributions are:
\begin{enumerate}[leftmargin=*]
    \item A generalization of SHAP to the structure of human motion, MotionSHAPE
    combining manifold-constrained perturbations with cohesive
    spatiotemporal masking, so that every reference input remains a
    plausible human movement. We instantiate MotionSHAP with two conditional generative models, a conditional VAE (MotionSHAP-VAE) and conditional flow matching (MotionSHAP-Flow), publicly available (upon publication of this manuscript).
    \item A suite of synthetic human-motion datasets with known ground-truth
    spatiotemporal feature attributions, enabling rigorous, quantitative evaluation of XAI
    methods for the first time in this domain. We show that classical attribution methods systematically miss these regions.
    \item An evaluation of XAI approaches on seven Parkinson’s disease gait-classification datasets from the CARE-PD benchmark. MotionSHAP systematically achieves the lowest Prediction Gap on Unimportant features (PGU)—the only metric that remains interpretable under the known tendency of self-evaluated deletion-style scores to favor overly disruptive imputations.
\end{enumerate}
Together, these contributions establish a principled foundation for
explainable AI in human motion analysis.